\documentclass[11pt]{article}
\usepackage[utf8]{inputenc}
\usepackage[T1]{fontenc}
\usepackage{geometry}
\usepackage{enumitem}
\geometry{a4paper, margin=1.8cm}

\begin{document}

\begin{center}
    {\LARGE \textbf{Gabarito – Cultura Digital:\\ Segurança Digital e Cidadania na Internet}}
\end{center}

\vspace{0.5cm}

\begin{enumerate}

\item \textbf{O que significa o termo “senha forte”?}  
\textbf{Resposta:} C) Uma senha com letras, números e símbolos.

\item \textbf{O que é phishing (pescaria virtual)?}  
\textbf{Resposta:} B) Um golpe que tenta enganar o usuário para roubar dados pessoais.

\item \textbf{Qual dessas atitudes aumenta a segurança nas redes sociais?}  
\textbf{Resposta:} B) Usar senhas diferentes para cada conta.

\item \textbf{O que é considerado um comportamento ético nas mídias digitais?}  
\textbf{Resposta:} C) Respeitar a privacidade e a opinião dos outros.

\item \textbf{Qual é a função principal do antivírus?}  
\textbf{Resposta:} C) Proteger o sistema contra programas maliciosos.

\item \textbf{O que é uma fake news?}  
\textbf{Resposta:} B) Uma notícia falsa criada para enganar ou manipular.

\item \textbf{Qual dessas ações ajuda a proteger seus dados pessoais na internet?}  
\textbf{Resposta:} C) Fazer login apenas em sites seguros (com “https”).

\item \textbf{O que significa o termo cyberbullying?}  
\textbf{Resposta:} C) A prática de ofender ou humilhar alguém por meio digital.

\item \textbf{Qual é uma característica das mídias digitais?}  
\textbf{Resposta:} C) Facilitam a interação e o compartilhamento de informações online.

\item \textbf{O que é backup?}  
\textbf{Resposta:} C) Uma cópia de segurança dos arquivos importantes.

\end{enumerate}

\end{document}

